% Written By Enoch Yu
% Downloaded from archive.enochyu.com

\documentclass{article}

\usepackage{amsmath}
\usepackage{amssymb}
\usepackage{amsthm}
\usepackage{mathtools}

\usepackage[margin=1in, headsep=15pt]{geometry}
\usepackage{lastpage}
\usepackage{fancyhdr}
\fancyhead[R]{Enoch Yu}
\fancyfoot[C]{Page \thepage \hspace{1pt} of \pageref{LastPage}}
\pagestyle{fancy}
\usepackage{parskip}

\usepackage{enumitem}
\usepackage[dvipsnames]{xcolor}
\usepackage[normalem]{ulem}
\usepackage{url}
\usepackage{hyperref}
\hypersetup{
    colorlinks=true,
    linkcolor=black,
    filecolor=magenta,      
    urlcolor=cyan,
    pdftitle={7.14.2025 AMC 10, 12 Free Class Slide},
    pdfpagemode=FullScreen,
}
\usepackage{tabularx}
\usepackage{arydshln}
\usepackage{float}
\usepackage{pdfpages}

\usepackage{tikz}
\usetikzlibrary{calc, angles, shapes.geometric}
\usepackage{tkz-euclide}
\usepackage{pgfplots}
\pgfplotsset{compat=1.9}
\usepackage[font=small,labelfont=bf]{caption}

\newtheorem{theorem}{Theorem}[section]
\newtheorem{lemma}[theorem]{Lemma}
\newtheorem*{lemma*}{Lemma}
\newtheorem{sublemma}{Lemma}[section]
\newtheorem{proposition}{Proposition}
\newtheorem{corollary}{Corollary}[theorem]
\newtheorem{example}{Example}[section]
\newtheorem*{example*}{Example}
\newenvironment{solution}{\begin{trivlist}\item[]{\bf Solution}}{\qed \end{trivlist}}


\title{7.14.2025 AMC 10, 12 Free Class Slide}
\author{Enoch Yu}
\date{14 July 2025}

\begin{document}
\section*{Problem 1.} % 2022 AMC 10A Problem 16 (30)
The roots of the polynomial $10x^3 - 39x^2 + 29x - 6$
are the height, length, and width of a rectangular box
(right rectangular prism). A new rectangular box is
formed by lengthening each edge of the original box by
$2$ units. What is the volume of the new box? (Source: AMC 10)

\newpage
\section*{Problem 2.} % 2023 AMC 10A Problem 21 (47)
Let $P(x)$ be the unique polynomial of minimal degree
with the following properties:
\begin{itemize}
    \item $P(x)$ has a leading coefficient $1$,
    \item $1$ is a root of $P(x)-1$,
    \item $2$ is a root of $P(x-2)$,
    \item $3$ is a root of $P(3x)$, and
    \item $4$ is a root of $4P(x)$.
\end{itemize}
The roots of $P(x)$ are integers, with one exception.
The root that is not an integer can be written as
$\frac{m}{n}$, where $m$ and $n$ are relatively prime
integers. What is $m+n$? (Source: AMC 10)

\newpage
\section*{Problem 3.} % 2017 AMC 10A Problem 24 (-7007)
For certain real numbers $a$, $b$, and $c$, the polynomial
\[
g(x) = x^3 + ax^2 + x + 10
\]
has three distinct roots, and each root of $g(x)$
is also a root of the polynomial
\[
f(x) = x^4 + x^3 + bx^2 + 100x + c.
\]
What is $f(1)$? (Source: AMC 10)
\end{document}

